% =============================================================================
% 1.6 RAG : génération augmentée par récupération
% =============================================================================
\section{Génération augmentée par récupération}
\label{sec:ch1-rag}

\subsection{Définition et principes fondamentaux}

La génération augmentée par récupération (\anglais{Retrieval-Augmented
Generation}, RAG) est une architecture introduite par Lewis et coll.\ (2020)
qui combine un moteur de recherche et un générateur de texte. Plutôt que de
s'en remettre uniquement aux connaissances paramétriques du modèle, le système
récupère des extraits pertinents dans une base externe, puis conditionne la
génération sur ces extraits.

\begin{btpdef}[Génération augmentée par récupération]
On appelle génération augmentée par récupération toute architecture dans
laquelle un modèle de langage produit une réponse \emph{conditionnellement} à
un ensemble d'extraits récupérés au moment de la requête, plutôt qu'à partir
de sa seule mémoire paramétrique.
\end{btpdef}

Cette architecture répond à trois faiblesses des modèles de langage utilisés
seuls~: l'\btpterm{hallucination} (production d'affirmations non fondées),
l'\btpterm{obsolescence} (connaissances figées à la date d'entraînement) et
l'\btpterm{absence de citation} (impossibilité de rattacher une phrase à une
source). Pour un dossier d'appel d'offres, ces trois faiblesses sont
rédhibitoires.

La figure~\ref{fig:principe-rag} en rappelle le principe.

\begin{figure}[htbp]
  \centering
  \resizebox{\linewidth}{!}{% Principe général de la génération augmentée par récupération
\begin{tikzpicture}[node distance=6mm and 11mm]

  \node[btpboxstrong, text width=2.5cm] (q) {Question\\utilisateur};
  \node[btpbox, text width=2.9cm, right=of q] (search)
    {Moteur de\\recherche\\[-1pt]\btplabel{lexical, sémantique, hybride}};
  \node[btpdb, right=of search, minimum width=2.3cm] (kb)
    {Base de\\connaissances};
  \node[btpbox, text width=2.9cm, below=16mm of search] (ctx)
    {Contexte\\[-1pt]\btplabel{extraits pertinents numérotés}};
  \node[btpbox, text width=2.5cm, left=of ctx] (prompt)
    {Invite\\[-1pt]\btplabel{question + contexte + consignes}};
  \node[btpboxalt, text width=2.9cm, right=of ctx] (llm)
    {Modèle de\\langage};
  \node[btpboxstrong, text width=2.7cm, below=14mm of llm] (rep)
    {Réponse\\avec citations};

  \draw[btparrow] (q) -- (search);
  \draw[btparrow] (search) -- (kb);
  \draw[btparrowlight] (kb.south) |- (ctx.east);
  \draw[btparrow] (ctx) -- (prompt);
  \draw[btparrow] (prompt.south) |- ($(rep.west)+(-0.7,0)$) -- (rep.west);
  \draw[btparrow] (ctx) -- (llm);
  \draw[btparrow] (llm) -- (rep);

  \node[btplabel, below=1mm of rep, text width=6cm]
    {Chaque affirmation est rattachable à un extrait source identifié};

\end{tikzpicture}
}
  \caption{Principe de la génération augmentée par récupération}
  \label{fig:principe-rag}
\end{figure}

\subsection{Les composantes d'un système RAG}

\subsubsection{Le récupérateur}

Le récupérateur identifie les extraits pertinents dans la base. Les approches
examinées à la section~\ref{sec:ch1-recherche} s'y appliquent directement,
résumées au tableau~\ref{tab:approches-retriever}.

\begin{btptable}{Approches de récupération dans un système RAG}
\label{tab:approches-retriever}
\begin{tabularx}{\linewidth}{@{}L{2.6cm} Y L{3.4cm}@{}}
\toprule
\rowcolor{btppale}
\btpth{Approche} & \btpth{Principe} & \btpth{Cas d'usage} \\
\midrule
Lexicale & Classement par correspondance de termes (BM25). & Identifiants, articles, normes. \\
Sémantique & Similarité vectorielle dans un espace de plongement. & Concepts, paraphrases. \\
Hybride & Fusion des deux classements (RRF). & Compromis général. \\
Structurée & Filtrage sur métadonnées (type de pièce, organisation). & Isolation et précision. \\
Graphe & Navigation le long de relations explicites. & Questions multi-sauts. \\
\bottomrule
\end{tabularx}
\end{btptable}

\subsubsection{Le ré-ordonnançeur}

Le ré-ordonnançeur affine le classement produit par le récupérateur. Un modèle
croisé évalue la pertinence d'un couple (question, extrait) plutôt qu'une
distance entre vecteurs calculés indépendamment. Ce gain de précision est
décisif~: le contexte transmis au générateur est borné, et chaque extrait
hors sujet occupe une place qui aurait pu revenir à une source utile.

\subsubsection{Le générateur}

Le générateur est un modèle de langage qui reçoit la question, les extraits
récupérés et un ensemble de consignes, puis produit une réponse. Trois
contraintes s'imposent dans le domaine des appels d'offres~:

\begin{itemize}
  \item la réponse doit être \btpterm{citable}~: chaque affirmation factuelle
        est rattachable à un extrait identifié~;
  \item la réponse doit être \btpterm{structurée}~: le livrable n'est pas un
        paragraphe libre, mais une section dont le plan est imposé~;
  \item la réponse doit \btpterm{s'abstenir} plutôt qu'inventer lorsque le
        contexte ne contient pas l'information demandée.
\end{itemize}

\begin{btpnote}[Abstention plutôt qu'invention]
Dans un système conversationnel généraliste, une réponse approximative est
souvent acceptable. Dans un dossier d'offre, elle est plus dangereuse qu'un
silence. L'abstention explicite --- « cette information n'apparaît pas dans
les sources consultées » --- est donc une sortie légitime, et non un échec
du générateur.
\end{btpnote}

\subsection{Les variantes du RAG}

\subsubsection{RAG naïf}

Dans sa forme élémentaire, chaque question déclenche une recherche unique
suivie d'une génération unique, sans mémoire des tours précédents ni
raisonnement multi-étapes. Cette architecture suffit aux questions
factuelles isolées. Elle échoue dès que la requête est anaphorique
(« et la pénalité~? ») ou qu'elle exige plusieurs recherches successives.

\subsubsection{RAG à mémoire}

Une mémoire des interactions précédentes permet de réécrire la requête
courante à la lumière du dialogue. L'anaphore « et la pénalité~? » devient
alors « clause de pénalité du marché dont le CCAP vient d'être discuté ».
Cette variante est utile dans l'espace de travail assisté, où l'opérateur
enchaîne des questions sur un même dossier.

\subsubsection{RAG hiérarchique}

Le découpage parent-enfant introduit deux granularités~: un extrait court,
indexé, sert à la recherche précise~; un extrait plus large, non indexé,
sert de contexte au générateur. Cette stratégie, détaillée au
chapitre~\ref{chap:methodologie}, répond au dilemme classique de la taille
d'extrait~: trop court, le contexte manque~; trop long, la recherche perd
en précision.

\subsubsection{RAG agentique}

Le modèle devient alors un agent qui décide des actions à entreprendre ---
rechercher, évaluer, rechercher à nouveau --- avant de répondre. Cette
variante est pertinente pour les questions complexes à plusieurs sauts. Elle
introduit toutefois une indétermination du chemin d'exécution, difficilement
compatible avec l'exigence de traçabilité d'un dossier engageant. Nous y
revenons à la section~\ref{sec:ch1-workflows}.

\subsubsection{RAG sur graphe}

Une variante récente, popularisée notamment par les travaux de Microsoft
Research (Edge et coll., 2024), utilise un graphe de connaissances comme
structure de récupération~: le document est d'abord projeté en entités et
relations, la recherche navigue dans ce graphe, et le contexte assemblé
reflète des chemins plutôt que des extraits isolés. Les avantages ---
raisonnement multi-sauts, réponses plus relationnelles --- s'accompagnent
d'un coût d'indexation élevé et d'une dépendance à la qualité de
l'extraction. Nous la situons comme une perspective de seconde phase.

\subsection{Évaluation des systèmes RAG}

L'évaluation d'un système RAG est double, car une défaillance peut naître
aussi bien de la récupération que de la génération.

\begin{btptable}{Métriques d'évaluation d'un système RAG}
\label{tab:metriques-rag}
\begin{tabularx}{\linewidth}{@{}L{3.2cm} Y L{3.4cm}@{}}
\toprule
\rowcolor{btppale}
\btpth{Métrique} & \btpth{Définition} & \btpth{Cible} \\
\midrule
Rappel à $K$ & Proportion des extraits pertinents présents dans les $K$ premiers. & Récupération \\
Précision à $K$ & Proportion d'extraits pertinents parmi les $K$ premiers. & Récupération \\
Fidélité & Degré auquel la réponse s'appuie exclusivement sur le contexte fourni. & Génération \\
Ancrage & Proportion d'affirmations rattachables à une source vérifiable. & Génération \\
Exhaustivité & Couverture des exigences ou des aspects de la question. & Génération \\
Exactitude & Conformité factuelle au corpus et au domaine. & Génération \\
\bottomrule
\end{tabularx}
\end{btptable}

Des protocoles automatisés, dont RAGAS (Es et coll., 2024), proposent
d'estimer la fidélité et l'ancrage par un modèle juge. Ils accélèrent
l'itération mais ne dispensent pas d'une évaluation par des experts du
domaine, seule à même de juger la recevabilité d'une section d'offre.
